\begin{solapartat}
De forma esquemàtica podem representar les situacions de ressonància en les gràfiques següents:

\noindent\begin{minipage}[b]{0.45\linewidth}
\figura[1]{sol_fig1.png}
\end{minipage}\ \punts{0,2}

La relació entre la longitud d'un tub sonor i la longitud d'ona en condició de ressonància és:
\[ L_n = \frac{\lambda}{4}(1 + 2n)\ \text{\punts{0,2}} \Rightarrow \]
\begin{align*}
L_n - L_{n-1} &= \frac{\lambda}{4}(1+2n) - \frac{\lambda}{4}\,(1 + 2(n-1)) = \frac{\lambda}{2}\ \text{\punts{0,2}}\\
&= 0,57\ \mathrm{m} - 0,19\ \mathrm{m} = 0,38\ \mathrm{m} \Rightarrow \lambda = 0,76\ \mathrm{m}\ \text{\punts{0,2}}
\end{align*}
La velocitat del so serà:
\[ v_\textit{so} = \lambda\nu = 334\ \mathrm{m/s}\ \text{\punts{0,2}} \]
\end{solapartat}

\begin{solapartat}
La intensitat de so rebuda serà:
\[ I = \frac{P_T}{4\pi d^2}\ \text{\punts{0,3}} = 8,84\cdot 10^{-5}\ \mathrm{W/m^2}\ \text{\punts{0,2}} \]

Per tant el nivell de só en dB serà:
\[ \beta = 10\log\frac{I}{I_0} = \text{\punts{0,3}} = 79\ \mathrm{dB}\ \text{\punts{0,2}} \]
\end{solapartat}
